\title{Large Language Models Can Automatically Engineer Features for Few-Shot Tabular Learning}

\begin{abstract}
Large Language Models (LLMs) possess exceptional capabilities in addressing complex and novel reasoning challenges, presenting significant promise for tabular learning—a core component of numerous practical applications. In this work, we introduce \textsf{FeatLLM}, a new in-context learning framework that harnesses LLMs as feature engineers, generating representations of the input dataset that are specifically tailored to enhance tabular prediction tasks. These synthesized features are leveraged by straightforward downstream machine learning algorithms, such as linear regression, achieving robust performance in few-shot learning settings. Importantly, the \textsf{FeatLLM} framework requires only the simple predictive model and the constructed features for inference, circumventing the need for LLM access at prediction time for individual data points. Unlike prevailing LLM-based strategies, \textsf{FeatLLM} avoids repeated LLM queries for each test instance and only necessitates API-level interaction, thereby alleviating prompt size constraints. Empirical evaluations across a wide selection of tabular benchmarks from diverse sectors show that \textsf{FeatLLM} produces high-caliber rule sets, delivering superior accuracy—surpassing existing methods like TabLLM and STUNT by a margin of 10\% on average.
\end{abstract}

\section{Introduction}
Large Language Models (LLMs) have recently demonstrated striking proficiency in generating relevant responses to novel, unseen tasks without the need for explicit task-level fine-tuning~\cite{creswell2022selection,imani2023mathprompter,taylor2022galactica}. By being trained on vast and diverse textual resources, these models internalize a wealth of prior knowledge, enabling them to compensate for the absence of specialized domain expertise in multiple contexts~\cite{Genomes2021,singhal2023large,wilcox2003role}. This capability facilitates the automation of tasks across varied domains, such as dialogue interpretation~\cite{finch2023leveraging} and automated program correction~\cite{fan2023automated}. Notably, when provided with descriptions of the task and a small set of example instances via prompting, LLMs are able to capture task-relevant context and direct attention towards key features and exemplars~\cite{brown2020language,zhao2023survey}. Such behavior signals that LLMs are adept at data analysis and possess significant potential as automated feature engineers.

Recent advancements have broadened the application of LLMs’ reasoning and background knowledge to tabular learning problems. For example, knowledge encapsulating relationships like the increased susceptibility of older populations to particular diseases can inform and enhance disease prediction models. In this paradigm, recent methodologies articulate both tasks and features in natural language, convert tabular data into serialized text, and input these representations into an LLM~\cite{dinh2022lift,wang2023anypredict}. Such approaches frequently adopt in-context learning~\cite{nam2023semi} or leverage parameter-efficient adaptation techniques~\cite{dinh2022lift,hegselmann2023tabllm}. The integration of LLM-driven prior knowledge has been empirically validated to surpass traditional tabular learning baselines, particularly in scenarios where only limited data samples are available~\cite{hegselmann2023tabllm}.

Existing tabular learning techniques based on LLMs face several notable drawbacks. In the case of end-to-end prediction, each individual data point necessitates a separate LLM inference, incurring significant computational costs. Moreover, achieving strong predictive performance often relies on fine-tuning the LLMs, which is typically not practical for models lacking public training infrastructure or available tuning APIs. This limitation is particularly pronounced with state-of-the-art LLMs that restrict users to inference APIs only~\cite{anil2023palm,openai2023gpt4}. Another major issue concerns prompt length: as the feature set of a tabular dataset grows, prompt serialization inflates accordingly, rendering many current methods ineffective or degrading their predictive abilities when long prompts are involved~\cite{liu2023lost}. Collectively, these issues create substantial obstacles to deploying LLM-based models for tabular tasks in practice.

To address these issues, our proposed method departs from the direct use of LLMs for final prediction, instead leveraging them for feature engineering. This reframing enables the LLM to more effectively capitalize on its existing knowledge. We present \textsf{FeatLLM}, an innovative in-context learning paradigm designed to extract the "criteria" driving LLM prediction behavior. For example, in a medical diagnosis scenario, the LLM can be prompted to explicitly formulate logical rules or criteria identifying which feature patterns correspond to a disease diagnosis. These discovered criteria are then transformed into engineered features that substitute for original input variables in downstream modeling tasks. In this workflow, the LLM functions primarily as a feature constructor, facilitating the use of straightforward models at inference, thus reducing prediction latency. Importantly, by leveraging in-context learning via prompt demonstrations, one can generate new features without engaging in any additional training. To tackle issues posed by large prompt sizes, we incorporate ensemble learning strategies such as feature bagging~\cite{breiman2001random}, thereby improving scalability to datasets with extensive feature sets.

\textsf{FeatLLM} decomposes the process of engineering novel features from input prompts into two main phases: (1) comprehending the underlying problem and deducing connections between features and target variables; and (2) using the insights from the first phase along with a small set of illustrative examples to develop concrete rules that separate classes. This organized reasoning framework guides LLMs to disregard irrelevant features during rule discovery. The generated rules are subsequently executed on the dataset (interpreted as program instructions), converting the raw data into binary indicators that signify whether each sample satisfies a specific rule. These newly formed binary features serve as inputs to a simple predictive model that estimates class probabilities. To bolster robustness, an ensemble strategy is used, which involves iteratively generating features while applying feature bagging. By tuning the count of sampled features and instances during bagging, the issue of oversized prompts is alleviated. \textsf{FeatLLM} requires no model retraining and maintains low inference overhead, supporting its deployment across diverse real-world scenarios.

Our assessment of \textsf{FeatLLM} spans 13 tabular datasets within low-sample settings, highlighting its consistently strong and stable results. We demonstrate that LLMs adeptly leverage both external prior knowledge and patterns identified from the dataset, resulting in effective feature extraction. Across a variety of evaluation conditions, our approach surpasses current state-of-the-art few-shot learning methods. The source code is available via an anonymized GitHub repository at~\texttt{https://github.com/Sungwon-Han/FeatLLM}.

\section{Related Work}

\subsection{Few-Shot Learning with Tabular Data} Few-shot learning has seen significant progress, enabling the development of broadly applicable data representations with only limited labeled examples~\cite{chen2018closer}. While research in this field began with a focus on images, its applicability has since expanded to encompass a variety of modalities, such as text, audio, and most recently, tabular data~\cite{majumder2022few,nam2023stunt,schick2021exploiting}. A key challenge when working with few labeled instances is the heightened risk of models capturing coincidental or misleading correlations~\cite{bartlett2021deep}. To address this, recent research has explored the use of external data sources or infused models with domain-relevant inductive biases through prior knowledge. Techniques in this vein include harnessing abundant unlabeled datasets and data augmentation for semi-supervised training~\cite{ucar2021subtab,yoon2020vime}, as well as pursuing effective model initialization using unsupervised, task-independent pre-training~\cite{somepalli2021saint}. Common unsupervised methods involve partially masking or perturbing data and relying on reconstruction-based loss functions~\cite{arik2021tabnet,majmundar2022met}, or using data augmentation to facilitate contrastive learning approaches~\cite{bahri2022scarf,chen2020simple}. Nonetheless, due to the intrinsic heterogeneity of tabular data, developing generic augmentation strategies remains a significant obstacle, with limited effectiveness observed in few-shot applications~\cite{nam2023stunt}.

More recent studies have advocated for model pre-training on a broad and diverse collection of tabular datasets—often unrelated to the eventual downstream task—before employing transfer learning techniques for adaptation~\cite{zhang2023generative,levin2022transfer,zhu2023xtab}. TransTab~\cite{wang2022transtab}, for example, utilizes transformers to extract semantic relationships among columns, making use of column names as well as table entries. Such cross-table, column-level knowledge has enhanced the capabilities of models to undertake zero-shot or few-shot reasoning on new tabular tasks. Similarly, TabPFN~\cite{hollmann2022tabpfn} synthesizes datasets with diverse distributions for the pre-training of a Bayesian neural network. Once pre-trained, the model performs inference by receiving both training and test data as input, without requiring additional task-specific training. The domain knowledge obtained from these varied sources has not only improved upon classic tree-based algorithms, but also yielded tangible gains under few-shot scenarios.

\subsection{Language-Interfaced Tabular Learning} An alternative and increasingly influential direction for injecting prior knowledge involves leveraging large language models (LLMs)~\cite{anil2023palm,openai2023gpt4}. Thanks to their exposure to vast and varied textual data during pre-training, LLMs display strong generalization abilities, making them effective tools across multiple problem domains~\cite{brown2020language}. Recent work has focused on tapping into the implicit knowledge within LLMs to address tabular tasks, often by reformatting tabular data into textual representations and incorporating these within LLM input prompts~\cite{dinh2022lift,hegselmann2023tabllm,wang2023anypredict}. By embedding the tables, along with contextual details such as feature and task descriptions, into prompts, models are able to infer the relationships between variables and tasks. The field has progressed to include methods that fine-tune these models with efficient parameter adaptation schemes, such as LoRA~\cite{hu2021lora} or IA3~\cite{liu2022few}, as well as strategies that rely solely on in-context learning—providing several task demonstrations within prompts without any gradient-based updating of model parameters~\cite{dinh2022lift,hegselmann2023tabllm,nam2023semi,slack2023tablet}.

Although prior approaches have shown promise in enhancing few-shot learning for tabular data, their reliance on routing all inferences through an LLM poses practical difficulties, especially in industrial contexts. In this work, we propose to move away from the traditional black-box paradigm, which involves processing individual sample predictions end-to-end with an LLM. Rather, we instruct the LLM to infer the discriminative rules or conditions specific to each class. \textsf{FeatLLM} subsequently transforms these derived rules into executable program code, generating novel features that allow predictions to be made independently of the LLM. This approach leverages in-context learning to distill rules based on both existing knowledge and the information present in a small number of labeled samples.

\section{Methods}
\textsf{FeatLLM} aims to derive the "rules"—that is, the decision criteria related to each target class—using only the provided few-shot examples, as shown in Figure~\ref{fig:main_model}. The rules are formulated by the LLM and then converted into binary features for each input instance, flagging whether individual rules are satisfied. These binary indicators are combined through a dot product with non-negative, learnable parameters to compute class probabilities. Through training, the model discerns the contribution of each extracted feature to the final prediction, as detailed in Section~\ref{Sec:training}. This framework incorporates repeated sampling using bagging and aggregates predictions using an ensemble strategy. The following sections elaborate on the problem formulation and detail the individual steps of the methodology.

\vspace*{-0.12in}
\paragraph{Problem formulation.} Consider a tabular dataset containing $N$ annotated instances, denoted by $\mathcal{D} = \{(\mathbf{x}^i, \mathbf{y}^i)\}_{i=1}^{N}$. Each instance $\mathbf{x}^i$ is a vector in $\mathbb{R}^d$, corresponding to $d$ distinct features. These features are labeled with descriptive names in natural language, represented as $F = \{f_j\}_{j=1}^{d}$, and are each accompanied by separate definitions. In the context of classification tasks, $\mathbf{y}^i$ denotes a category drawn from a predefined set. For $k$-shot learning scenarios, the model is trained on a randomly chosen subset of $k$ ($k < N$) labeled instances from the full dataset.

\subsection{Prompt Design for Extracting Rules}
\label{Sec:prompt_design}
To facilitate rule extraction by the LLM through more precise reasoning processes, we structure the prompt to encourage problem solving that aligns with expert human approaches to tabular learning. The prompt is organized into three primary components, illustrated in Fig.~\ref{fig:prompt_for_rule} and further detailed in Appendix~\ref{sec:example_rule}:

\vspace*{-0.12in}
\paragraph{Basic information description.} This section supplies the key context required to address the task (refer to the orange-highlighted text in Fig.~\ref{fig:prompt_for_rule}). It comprises a question-formulated description of the objective with the associated label set, definitions of the tabular features, and concrete examples. The task prompt poses the problem as a direct question (e.g., ``Does this patient's myocardial infarction complications data indicate chronic heart failure? Yes or no?"). Additional instances are listed in Appendix~\ref{sec:task_desc}. Each feature's description specifies its data type and, if relevant, a detailed definition. For features with categorical values, illustrative categories can be given. Example demonstrations serialize a small selection of labeled training samples into textual format, using ground-truth annotations. Serialization mirrors conventions established in prior works~\cite{dinh2022lift,hegselmann2023tabllm}:

\begin{align}
&\text{Serialize}(\mathbf{x}^i,\mathbf{y}^i, F) = \nonumber \\ 
&``\ f_1\ \text{is}\ \mathbf{x}^i_1.\ ... \ f_d\  \text{is}\  \mathbf{x}^i_d.\ \ \text{Answer:}\  \mathbf{y}^i\ ”
\end{align}

\vspace*{-0.12in}
\paragraph{Reasoning instruction.} Our objective is to improve the reasoning capabilities of the LLM by delivering explicit reasoning instructions (refer to the blue highlights in Fig.~\ref{fig:prompt_for_rule}). The reasoning instruction opens with a preliminary statement, echoing the style of the chain-of-thought methodology~\cite{wei2022chain}. Subsequently, the LLM’s inferential process is partitioned into two sequential phases. Initially, the LLM is prompted to hypothesize the underlying causal relationships or propensities between the various features and the defined task, leveraging broad knowledge or common sense instead of example demonstrations. This approach is intended to maximize the use of the LLM’s existing understanding relevant to the task. The following phase involves the LLM synthesizing both its initial causal insights and the presented demonstrations to generate classification rules for each target class. This bifurcated process serves to steer the model away from depending on coincidental associations found in irrelevant attributes, encouraging attention to the most pertinent features. To avoid the issues of underfitting through an insufficient number of rules or unwieldy prompt sizes due to an excess, we specify an exact quantity of rules to be formulated for each class (set at 10)\footnote{See Section~\ref{sec:analysis} for further discussion regarding rule count selection.}. Moreover, during rule generation, the LLM is instructed to utilize the feature types as provided in the foundational information segment, minimizing the risk of generating incompatible rules.

\vspace*{-0.12in}
\paragraph{Response instruction.} To streamline both the interpretation and utilization of generated rules, we instruct the LLM to adhere to a specified structure for its output responses (see yellow highlights in Fig.~\ref{fig:prompt_for_rule}). The prompt details the exact formatting expectations for each class of answers (i.e., response formatting) as well as the preferred layout for each individual rule (i.e., rule structure). By following these directives, the LLM is equipped to select and apply the correct format aligned with the specific feature types involved. In the absence of detailed constraints, the LLM may also combine rules by leveraging logical connectors such as AND or OR. A sample output generated under these instructions is presented in Appendix~\ref{sec:outcome-rule}.

\subsection{Inferring Class Likelihood via Rules}
\label{Sec:training}

\paragraph{Parsing rules for feature generation.} The previously derived rules are employed to construct new binary features for each class, where these features indicate whether a given sample fulfills the respective class's rules (taking a value of either '0' or '1'). Such features serve as predictors for estimating the probability that a sample pertains to each class. Because these rules, generated by the LLM, are expressed in natural language, there is a necessity to design a dedicated parsing approach that can systematically transform them into binary features. While formatting constraints are enforced during rule generation to streamline subsequent parsing, LLM outputs often display unanticipated syntactic variations~\cite{kovavc2023large}. For example, the LLM might replace symbolic representations like ``$>$'' or ``$<$'' with their verbal forms such as ``greater than'' or ``less than,'' introducing inconsistencies that complicate the parsing process.

To manage the complexities arising from noisy natural language text, rather than writing intricate parsing code, we capitalize on the LLM's semantic reasoning capabilities. The LLM excels at discerning underlying meanings despite fluctuations in syntax, and it is adept at producing concise code when properly prompted, due to its exposure to programming data. To exploit these capabilities, we construct prompts that incorporate the function name, detailed input and output specifications, and the rules to be parsed, and then present these prompts to the LLM. Illustrative examples of such prompts and their resultant outputs are provided in Figures~\ref{fig:prompt_for_parsing} and~\ref{sec:example_parsing}-\ref{sec:example_parse_outcome} in the Appendix. The generated parsing code is executed using Python's \texttt{exec()} mechanism to transform the raw data accordingly.

\vspace*{-0.12in}
\paragraph{Inferring class likelihood.} With the acquisition of binary features derived from class-specific rules, a basic approach to assess a sample’s likelihood for each class involves tallying the number of satisfied rules for each class—essentially summing the class's binary features. Nevertheless, because not all rules contribute equally, it becomes necessary to infer their relative importance through supervised learning using training data. We pursue this by fitting a traditional machine learning (ML) model to the binary features associated with each class, thereby quantifying feature weights. A typical choice is a linear model with no intercept term. Specifically, let $\mathbf{z}_k^i \in \{0, 1\}^R$ represent the binary feature vector for sample $\mathbf{x}^i$ under class $k$, where $R$ denotes the count of rules for each class and $c$ is the total number of classes. The predicted class probabilities $\mathbf{p}^i$ are then determined as follows:

\begin{align}
\text{logit}_k^i &= \max(\mathbf{w}_k, 0) \cdot \mathbf{z}_k^i, \nonumber \\
\mathbf{p}^i &= \text{Softmax}({[\text{logit}_{1}^i, ..., \text{logit}_{c}^i]}). \label{eq:infer_prob}
\end{align}

In this context, $\mathbf{w}_k$ denote the learnable weights associated with the linear model, encapsulating the relative significance assigned to each feature. By constraining these weights to reside within a positive subspace, we guarantee that the features produced by the LLM cannot negatively impact class probabilities during model training. Subsequently, the logit output undergoes a transformation into class probabilities via the softmax operation. Training of the weights $\mathbf{w}_k$ is accomplished through the minimization of cross-entropy loss, using a limited set of annotated examples under few-shot conditions.

\vspace*{-0.12in}
\paragraph{Ensembling with bagging.} The process is iterated multiple times, yielding a diverse set of inference models whose outputs are integrated using ensemble techniques for the ultimate prediction. Specifically, given a set of weights trained over $T$ independent runs, $\{\mathbf{w}[t]\}_{t=1}^T$, class probabilities $\mathbf{p}^i$ are estimated for each model and then averaged in accordance with Eq.~\ref{eq:infer_prob}.

Randomness is injected into the rule generation process for ensembling by setting a high temperature during LLM sampling. Furthermore, we shuffle the sequence of few-shot demonstrations because prior work has noted that LLM output can be sensitive to their ordering~\cite{lu2022fantastically}\footnote{Appendix~\ref{sec:diversity} provides a discussion on rule diversity.}. To further exploit diversity and improve predictive stability, bagging is employed to randomly choose subsets of features or data instances for each trial—an approach that is increasingly beneficial as the volume of data or feature set grows. This ensemble methodology confers two notable strengths. First, should some trials yield faulty rules, correct rules identified in other iterations compensate, thereby strengthening the framework's resilience. This benefit is related to the LLM's capacity for self-consistency~\cite{wang2022self}; rules repeatedly surfaced across different trials are statistically more likely to be correct. Second, the use of ensembles mitigates constraints imposed by limited prompt sizes in LLMs: bagging naturally restricts the number of features or samples per trial, supporting reliable operation even when the data's size would otherwise exceed prompt limits. Although a single collection of rules may overlook associations involving some features, assembling many weak learners from various trials can effectively offset issues arising from partial feature or instance representation.

\section{Experiments}
We assess the effectiveness of \textsf{FeatLLM} across a diverse collection of tabular datasets and perform multiple analyses to better understand the mechanics of our approach. Supplementary experiments—including those that test different LLMs and qualitative investigations—are presented in the Appendix due to space limitations.

\subsection{Performance Evaluation}
\paragraph{Datasets.}
Our empirical study involves 13 datasets for tasks involving binary or multi-class classification: (1) Adult~\cite{asuncion2007uci} aims to predict whether a person's yearly income exceeds \$50,000; (2) Bank~\cite{moro2014data} predicts the likelihood of a customer subscribing to a term deposit; (3) Blood~\cite{yeh2009knowledge} forecasts whether a blood donor will make additional donations; (4) Car~\cite{kadra2021well} assesses car quality; (5) Communities~\cite{misc_communities_and_crime_183} estimates regional crime rates; (6) Credit-g~\cite{kadra2021well} evaluates whether an individual is a favorable or risky credit candidate; (7) Diabetes\footnote{\texttt{kaggle.com/datasets/uciml/pima-indians-diabetes-database}} determines diabetes status; (8) Heart\footnote{\texttt{kaggle.com/datasets/fedesoriano/heart-failure-prediction}} predicts coronary artery disease incidence; and (9) Myocardial~\cite{misc_myocardial_infarction_complications_579} checks for chronic heart failure. 

In addition, our benchmark comprises newer datasets and synthetic ones that are rarely explored and were omitted from LLMs’ pre-training corpora: (10) Cultivars, used for predicting the expected grain yield of a soybean cultivar\footnote{\texttt{archive.ics.uci.edu/dataset/913}}; (11) NHANES, for inferring a person’s age group from medical records\footnote{\texttt{archive.ics.uci.edu/dataset/887}}; (12) Sequence-type, classifying sequences as either arithmetic, geometric, Fibonacci, or Collatz; and (13) Solution-mix, assessing whether the mixture concentration from four solutions surpasses 0.5. The first two datasets have been newly released (post-September 2023) to the UCI Machine Learning Repository, ensuring their exclusion from the pre-training data of the LLM API (gpt-3.5-turbo-0613) employed in our approach. The last two, being synthetically constructed, could not be recalled or memorized by the LLM API, thus enabling a fair evaluation.

The dataset selection features a broad spectrum in both sample sizes and attribute complexities (see Table~\ref{Tab:basic-desc}). Each collection includes explicit attribute names and explanations. Some, such as `Communities' and `Myocardial,' present a challenge with their extensive attribute counts (exceeding 100), particularly considering the finite context length constraints inherent in LLM architectures.

\vspace*{-0.12in}
\paragraph{Baselines.}
We compare our approach against nine alternative methods. The first set includes three standard tabular data learners: (1) Logistic regression (LogReg), (2) XGBoost~\cite{chen2016xgboost}, and (3) RandomForest~\cite{ho1995random}. The subsequent two—(4) SCARF~\cite{bahri2022scarf} and (5) STUNT~\cite{nam2023stunt}—utilize unlabeled data during training, although acquiring substantial amounts of such data may be impractical in realistic settings. A more contemporary technique, (6) TabPFN~\cite{hollmann2022tabpfn}, pre-trains via synthetically generated datasets sampled from a variety of distributions. The remaining baselines employ LLMs: (7) In-context learning (In-context)~\cite{wei2022emergent}, (8) TABLET~\cite{slack2023tablet}, and (9) TabLLM~\cite{hegselmann2023tabllm}. The in-context approach embeds a handful of training examples into the model prompt without model fine-tuning. TABLET augments the prompts with interpretative hints using rule sets and prototype-based summaries via an external classifier to improve inference. Finally, TabLLM leverages efficient parameter tuning strategies, such as IA3~\cite{liu2022few}, to optimize LLM performance when limited annotated samples are available.

\vspace*{-0.12in}
\paragraph{Implementation details.}
\textsf{FeatLLM} employs GPT-3.5 as its default LLM backbone, though the framework is built to be flexible regarding the underlying LLM used (refer to Appendix~\ref{sec:backbones} for evaluations with alternative backbones). For LLM inference, we fix the temperature parameter at 0.5 and retain the top-p setting at the default API value of 1. We configure the ensemble count and the number of extraction rules to 20 and 10, respectively. A comprehensive analysis of the influence of various hyper-parameters is provided in Figure~\ref{fig:hyperparameter2} and detailed in Appendix~\ref{sec:hyper-parameter}. The linear model is trained using the Adam optimizer with a learning rate of 0.01 for 200 epochs, utilizing $k$-fold cross-validation to select the best epoch. For the baseline experiments, GPT-3.5 facilitates in-context learning, while TabLLM leverages T0~\cite{sanh2022multitask}, as per the original setup. It is important to note that TabLLM permits only LLMs with publicly available checkpoints, thus excluding GPT-3.5. In the context of standard machine learning algorithms like LogReg, XGBoost, and RandomForest, hyper-parameters are identified through a Grid-search strategy, coupled with $k$-fold cross-validation. Here, $k$ is set to either 2 or 4, ensuring that the training set contains a minimum of one example for every class. Further implementation specifics are documented in Appendix~\ref{sec:baselines}.

\vspace*{-0.12in}
\paragraph{Results.}
Tables~\ref{tab:main1} and \ref{tab:main2} display comparative performance metrics across several datasets. Each experiment is conducted three times with different random splits of the training and test sets; we report both the mean and standard deviation for the AUC scores. Across all evaluated settings, \textsf{FeatLLM} consistently achieves either the highest or second-highest ranking among the compared baseline methods. Notably, our framework delivers performance on par with traditional tabular baselines requiring 32 shots, using only 4 shots (refer to Fig.~\ref{fig:compares}). While STUNT and TabPFN exhibited notable improvements on particular datasets, their aggregate advantage over classical machine learning models was limited. Additionally, LLM-based models like in-context learning, TABLET, and TabLLM were unable to perform inference on tabular datasets characterized by high-dimensional feature spaces. By integrating bagging and ensemble strategies, our framework maintained robust results even for datasets with a large number of features. It should be emphasized that incorporating our bagging and ensemble approach into current LLM-based models is feasible, though doing so would require duplicating the number of inference operations per sample, resulting in considerably higher computational overhead and thus rendering such approaches impractical.

\subsection{Analysis \& Discussion}
\label{sec:analysis}
\paragraph{Ablation study.} To assess the significance of key framework components, we carry out ablation studies examining: (1) \textsf{-Tuning}: the exclusion of the weight-tuning step in the linear model, using a simple sum of binary features for logit calculation instead; (2) \textsf{-Ensemble}: the removal of the ensemble step; (3) \textsf{-Description}: the omission of feature descriptions; and (4) \textsf{-Reasoning}: the removal of the initial reasoning step (Step-1) in the rule generation phase.

Table~\ref{tab:ablation} reports the mean change in AUC due to these ablations, averaged across all datasets. In all cases, omitting any component causes a reduction in performance. Notably, the positive impact of weight tuning becomes more pronounced as the shot count increases, indicating that with more training examples, the model can better estimate rule importance, enhancing the effectiveness of weight tuning. Conversely, the contributions of feature descriptions and reasoning instructions are especially substantial at lower shot counts, implying that leveraging LLMs' prior knowledge is critical when data is scarce. Finally, the adoption of the ensemble method yields consistent improvements in all tested scenarios.

\vspace*{-0.12in}
\paragraph{Training \& inference time comparison.} We assess and contrast the training and inference durations across various approaches. The experimental evaluation utilizes the Adult dataset, with the training set limited to 16 examples. All methods, aside from TabLLM, are executed on a single A100 GPU, while TabLLM employs four A100 GPUs to facilitate model parallelism. The detailed runtime statistics are reported in Table~\ref{tab:cost}. Importantly, \textsf{FeatLLM} demonstrates a relatively low inference latency, on par with standard tabular learning methods. By comparison, inference with other LLM-driven solutions—namely In-context learning, TABLET, and TabLLM—is significantly slower. When it comes to training, \textsf{FeatLLM} also shows comparable performance to the alternative LLM-based approaches; it only requires 30 API calls, an amount that neither substantially increases costs nor impedes scalability due to the feasibility of parallel execution.

\vspace*{-0.12in}
\paragraph{Quality of features generated by \textsf{FeatLLM}.}
To examine the effectiveness of the rule-based features generated by \textsf{FeatLLM} in real-world tasks, we carried out a straightforward empirical analysis. We quantified the informativeness of the generated features by computing the average AUROC between each feature and the target variable. We then calculated the proportion of features surpassing an AUROC of 0.5 for every dataset. The findings, summarized in Table~\ref{tab:quality}, indicate that a substantial share of the generated rules are strongly associated with the prediction target and contribute meaningfully to solving the main task (see ``Before tuning" column). Furthermore, as the linear model is fitted to the data, rules with marginal utility (those for which the assigned weights fall below a specified cutoff) are pruned automatically (refer to ``After tuning"). For these experiments, we set the cutoff value to 0.1 based on the distribution of learned weights.

\vspace*{-0.12in}
\paragraph{Effect of adding spurious correlations.} To evaluate the robustness of different approaches in filtering out extraneous spurious correlations under a low-shot setting, we designed an experiment. Specifically, we selected the Heart dataset and randomly appended to it columns drawn from the Adult dataset that are not pertinent to the Heart dataset’s predictive task. We then monitored model performance as we systematically increased the number of irrelevant features. For equitable assessment, each method was restricted to 8 labeled examples, with no access to any unlabeled data—thereby omitting approaches like SCARF and STUNT from the scope of comparison.

Figure~\ref{fig:spurious} displays how the models' performance responds to the inclusion of spurious correlations. Notably, \textsf{FeatLLM} experienced the smallest drop in accuracy as more irrelevant features were added, outperforming baseline approaches. This robustness likely stems from the framework’s reliance on prior knowledge to determine feature-task associations when formulating rules, thereby mitigating the incorporation of rules centered on irrelevant columns. As a result, the fraction of rules generated from noisy features was observed to be just 1-2\% of the overall set. However, it is important to note that if unrelated but semantically similar features are introduced from related datasets, \textsf{FeatLLM} may still capture spurious patterns and misidentify causal links between features and the task—details of which are discussed further in Appendix~\ref{sec:spurious-appendix}.

\vspace*{-0.12in}
\paragraph{Hyper-parameter analysis}
\textsf{FeatLLM} operates with two principal hyper-parameters: the ensemble count, and the number of rules per class generated in each LLM query. To evaluate their influence on model effectiveness, we carried out a series of experiments. Our findings indicate that as the number of ensembles grows, performance also rises, yet these improvements plateau when the ensemble size surpasses roughly 20 (refer to Fig.~\ref{fig:appendix-hyperparameter1} in Appendix). Concerning the quantity of rules per class, we observed no statistically significant performance fluctuation. Nevertheless, considering the trade-off between prompt size and efficacy, we determined that extracting 10 rules per class achieves the most balanced outcome (see Fig.~\ref{fig:hyperparameter2}). We hypothesize that increasing rule numbers may raise the input's dimensionality, which enhances information content but also heightens overfitting risks, particularly in low-shot learning settings.

\section{Conclusion}
This work introduces \textsf{FeatLLM}, establishing new benchmarks for few-shot tabular learning utilizing LLMs. Unlike standard approaches that primarily rely on LLMs for individual predictions, \textsf{FeatLLM} leverages LLMs to articulate decision criteria, akin to a feature engineering process. By employing rule induction and aggregating outcomes using bagging, \textsf{FeatLLM} maintains low inference requirements, functions solely through API access without additional training, and remains robust to constraints in feature dimensionality. These design choices enable \textsf{FeatLLM} to deliver strong predictive performance, rendering it an effective and data-efficient solution for deploying LLMs on practical tabular datasets.

At present, \textsf{FeatLLM} is oriented toward low-shot learning situations. Looking forward, we aim to extend its framework to support the creation of new features for larger datasets, investigate alternative feature modalities beyond rule-based extraction, and enhance interpretability—thereby broadening its applicability to more diverse machine learning tasks.

\section{Broader Impact}
The \textsf{FeatLLM} framework seeks to harness LLMs' inherent prior knowledge to advance performance in tabular learning tasks, surpassing the limits of conventional supervised methods. Our approach specifically addresses data efficiency by countering overfitting stemming from limited labeled data, using LLM-informed augmentation. This makes \textsf{FeatLLM} particularly advantageous in applied domains (e.g., finance, healthcare) where labeling data is expensive or difficult and where LLMs—having been pre-trained on relevant domain-specific corpora—can significantly augment knowledge transfer. A critical avenue for future exploration, crucial for real-world adoption, involves systematically studying the societal implications of using LLM-derived prior knowledge, especially with regard to embedded biases or misinformation, since such priors may disproportionately shape model outputs relative to the available labeled data.

\end{document}